\documentclass[12pt]{article}
\usepackage[utf8]{inputenc}
\usepackage{mathpazo}
\usepackage{graphicx}
\usepackage[a4paper,width=150mm,top=25mm,bottom=25mm]{geometry}

\title{
{\includegraphics[width=0.5\textwidth]{logo.jpg}}\\
{\large \textbf{Bachelor of Computer Applications}}\\
Title of the project}
\author{Student name \\ Supervised by:Mr./Ms./Dr. Guide name}
\date{\today}

\begin{document}
\maketitle

\newpage

\tableofcontents

\newpage
\section{Abstract}

\section{Introduction}
\subsection{Background}
Please start from very big domain for your problem then focus on some area inside this domain that match your interest.
\subsection{Motivation}
Discus as much as market needs and academic needs for your project. It is expected that in this section you will get about 10~15 general technical CS challenges. Use surveys for showing the market need of your project.
Why this system is to be developed.
Your goals behind working and achieving your project
\subsection{Problem Definitions}
Please put a focus on 1~2 challenges that this project aims to solve and state them very clear as your formal problem statement.

\section{Project Description}
Show down with a figure the proposed system.
\subsection{Scope}
The scope of your project.
\subsection{Project Overview}
Some extra details for your system implementations.


\section{Similar System Information}
List down at least 4 papers from ACM and IEEE for similar work experience in the domain of your problem. You can add 5 papers or more from other sources (Springer, Elsevier, Website ..etc)

Be sure that each paper you list include the following points

\begin{enumerate}
\item Motivation of this work (Why the researchers do it)
\item The main problem statement of the work.
\item How the researchers contributed to solve the problem
\item What main results the researchers reach.
\item How do you think this paper you read is important for you. 
\end{enumerate}

\subsection{Similar System Description}
A system proposed by Rehm at al. \cite{Rehm} has been developed to ... Also the work was discussed in [\cite{Weiser},\cite{abc123}].


\bibliographystyle{IEEEtranS}
\bibliography{myref}
\end{document}
